% Stacking order follows layup.json: ply 1 bottom surface; ply 20 top surface.
\begin{table}[htbp]
\centering
\caption{Specified 20-ply laminate, stacked from bottom to top. Zero degrees is the beam longitudinal direction.}
\label{tab:layup}
\begin{tabular}{rllrr}
\toprule
Ply & Material & Orientation & Cured thickness (mm) & Fiber areal mass (g/m$^2$) \\
\midrule
1 & CF 2×2 twill, 200 gsm & 0^\circ & 0.20 & 200 \\
2 & UD E-glass/epoxy, 100 gsm & 0^\circ & 0.15 & 100 \\
3 & CF 2×2 twill, 200 gsm & 0^\circ & 0.20 & 200 \\
4 & UD E-glass/epoxy, 100 gsm & 0^\circ & 0.15 & 100 \\
5 & CF 2×2 twill, 200 gsm & 0^\circ & 0.20 & 200 \\
6 & UD E-glass/epoxy, 100 gsm & 0^\circ & 0.15 & 100 \\
7 & CF 2×2 twill, 200 gsm & 0^\circ & 0.20 & 200 \\
8 & UD E-glass/epoxy, 100 gsm & 0^\circ & 0.15 & 100 \\
9 & CF 2×2 twill, 200 gsm & 0^\circ & 0.20 & 200 \\
10 & UD E-glass/epoxy, 100 gsm & 0^\circ & 0.15 & 100 \\
11 & CF 2×2 twill, 200 gsm & 0^\circ & 0.20 & 200 \\
12 & UD E-glass/epoxy, 100 gsm & 0^\circ & 0.15 & 100 \\
13 & CF 2×2 twill, 200 gsm & 0^\circ & 0.20 & 200 \\
14 & UD E-glass/epoxy, 100 gsm & 0^\circ & 0.15 & 100 \\
15 & CF 2×2 twill, 200 gsm & 0^\circ & 0.20 & 200 \\
16 & UD E-glass/epoxy, 100 gsm & 0^\circ & 0.15 & 100 \\
17 & CF 2×2 twill, 200 gsm & 0^\circ & 0.20 & 200 \\
18 & UD E-glass/epoxy, 100 gsm & 0^\circ & 0.15 & 100 \\
19 & CF 2×2 twill, 200 gsm & 0^\circ & 0.20 & 200 \\
20 & UD E-glass/epoxy, 100 gsm & 0^\circ & 0.15 & 100 \\
\bottomrule
\end{tabular}
\end{table}

The stack is $[\mathrm{CF}_{0}/\mathrm{GF}_{0}]_{10}$ from bottom to top. It
contains 10 plies of each material, for a total cured thickness of 3.50 mm and
specified dry-fiber areal mass of 3.0 kg/m$^2$. The sequence is unsymmetric
through the thickness, so extension--bending coupling is possible.

\begin{table}[htbp]
\centering
\caption{Representative orthotropic elastic inputs used by the layered-shell model. Moduli are in GPa.}
\label{tab:ply-properties}
\begin{tabular}{lrrrrrrrrrr}
\toprule
Material & $E_1$ & $E_2$ & $E_3$ & $\nu_{12}$ & $\nu_{23}$ & $\nu_{13}$ & $G_{12}$ & $G_{23}$ & $G_{13}$ & $\rho$ (kg/m$^3$) \\
\midrule
CF 2$\times$2 twill & 60.6 & 61.6 & 8.3 & 0.10 & 0.25 & 0.25 & 5.00 & 3.00 & 3.00 & 1724 \\
UD E-glass/epoxy & 38.6 & 8.27 & 8.27 & 0.26 & 0.30 & 0.26 & 4.14 & 3.18 & 4.14 & 1111 \\
\bottomrule
\end{tabular}
\end{table}

For carbon twill, the 0/90 tensile moduli are taken from a supplier's 2$\times$2
twill prepreg data; the reported 42\% resin content is used for a simple
cured-density estimate. The missing in-plane shear and transverse constants
are representative modeling assumptions. For glass, the input assumes
unidirectional E-glass/epoxy and uses representative published lamina
constants; the transverse constants are approximated with transverse symmetry.
The glass resin mass fraction is assumed to be 40\%. These are deliberately
approximate numerical inputs, not measured properties for the user's material
system. The estimated cured laminate areal mass is 5.115 kg/m$^2$.

The carbon values are based on the Gurit SE84LV 2$\times$2 twill prepreg
properties, with selected missing values taken from generic carbon-fabric
references \cite{guritSE84LV,performanceCarbonGeneric}. The glass in-plane
properties are representative UD E-glass/epoxy values \cite{glassUDLamina}.
